\documentclass{eleclab-report}

\setexperimentnumber{3}
\setexperimenttitle{直流电路定理的研究}

\begin{document}
\makelabtitle
\makelabheadertable

\LabMajorSection{实验目的}
\LabBlank{3.1\baselineskip}

\LabMajorSection{实验原理}
\begin{LabArabicList}
  \item 叠加性和齐次性
  \item 戴维南定理和诺顿定理
  \item 测量有源二端网络等效内阻
\end{LabArabicList}
\LabBlank{0.8\baselineskip}

\LabMajorSection{实验设备}
\LabBlank{2.8\baselineskip}

\LabMajorSection{实验内容}

\LabMinorSection{验证线性电路叠加定理及齐次性}
\LabPlaceholderBlock{电路图}{4.6\baselineskip}
\LabPlaceholderBlock{实验步骤}{2.0\baselineskip}
\LabPlaceholderBlock{测试设备}{2.0\baselineskip}
\LabItem{测试数据}
\LabSuperpositionTable
\LabDataAnalysisBlock{3.8\baselineskip}

\LabMinorSection{测量有源二端网络的等效内阻}
\LabPlaceholderBlock{电路图}{4.3\baselineskip}
\LabPlaceholderBlock{实验步骤}{2.0\baselineskip}
\LabPlaceholderBlock{测试设备}{2.0\baselineskip}
\LabItem{测试数据}
\LabMethod{1. 开路电压-短路电流法}
\LabOpenShortTable
\LabMethod{2. 伏安法}
\LabVoltAmpereTable
\LabMethod{3. 直接法}
\LabDirectMethodTable
\LabDataAnalysisBlock{3.4\baselineskip}

\LabMinorSection{验证戴维南定理}
\LabPlaceholderBlock{电路图}{4.2\baselineskip}
\LabPlaceholderBlock{实验步骤}{2.0\baselineskip}
\LabPlaceholderBlock{测试设备}{2.0\baselineskip}
\LabItem{测试数据}
\LabTheveninTable
\LabDataAnalysisBlock{4.0\baselineskip}

\LabMinorSection{验证诺顿定理}
\LabPlaceholderBlock{电路图}{4.2\baselineskip}
\LabPlaceholderBlock{实验步骤}{2.0\baselineskip}
\LabPlaceholderBlock{测试设备}{2.0\baselineskip}
\LabPlaceholderBlock{测试数据}{3.0\baselineskip}
\LabDataAnalysisBlock{3.0\baselineskip}

\LabMajorSection{思考题}
\labresetitems
\LabItem{在本实验中验证叠加原理时，若电压源内阻不能忽略，实验该如何进行？}
\LabBlank{1.5\baselineskip}
\LabItem{叠加原理的使用条件是什么？}
\LabBlank{1.5\baselineskip}
\LabItem{在本实验中验证诺顿定理时，如果忽略电流源内阻实验结果会怎样？}
\LabBlank{1.5\baselineskip}
\LabItem{组成戴维南等效电路时，电压源U'oc要用表5.2.2中用来测量UOC的万用表的同一个量程进行校准、电阻箱R'eq要用万用表的欧姆档进行校准，不能以电压源、电阻箱上的读数为准，为什么？}

\end{document}
